% Unite: papier 29; traduction francaise depuis la source allemande auditee.
\setcounter{footnote}{0}

\section*{29. Le théorème de finitude des invariants des groupes linéaires finis de caractéristique $p$}

\begin{center}
Nachr. v. d. Ges. d. Wiss. zu Göttingen 1926, p. 28--35
\end{center}

\begin{center}
Par Emmy Noether à Göttingen.\\[0.5em]
Présenté par R. Courant à la séance du 30 juillet 1926.
\end{center}

Je montre dans ce qui suit que le théorème bien connu de finitude des invariants des groupes finis de substitutions linéaires\footnote{Pour une démonstration élémentaire de finitude, voir E. Noether, Der Endlichkeitssatz der Invarianten endlicher Gruppen. Math. Ann. 77 (1916), p. 89--92.} reste valable lorsque, comme domaine des coefficients, on prend au lieu d'un corps de nombres ordinaire un corps quelconque -- donc en particulier un corps de caractéristique $p$.\footnote{Pour les notions de théorie des corps, voir E. Steinitz, Algebraische Theorie der Körper. J. f. M. 137 (1910), p. 167--309. Caractéristique au \S{} 4, corps des racines au \S{} 12, corps des fractions au \S{} 3.} Mais les démonstrations de finitude connues échouent dès que l'ordre du groupe est divisible par $p$; il faut faire intervenir des faits arithmétiques plus profonds, à savoir le critère suivant, connu jusqu'ici seulement pour la caractéristique nulle:

\textbf{Critère de finitude}\footnote{Comparer E. Noether, Algebraische und Differentialinvarianten. Jahresber. d. d. Math.-Ver. 32 (1923), p. 177--184, no 3, théorème 4, ainsi que la littérature qui y est indiquée. Dans la dernière phrase du critère il y est dit par erreur ``base de $\mathfrak S$ comme $\mathfrak R$-module'' (système fondamental des grandeurs entières) au lieu de ``comme module sur un certain sous-anneau de $\mathfrak R$''.}: Un domaine d’intégrité $\mathfrak S$ formé de polynômes -- plus généralement de fonctions rationnelles -- à coefficients dans un corps est fini si et seulement si $\mathfrak S$ contient un sous-anneau fini $\mathfrak R$ tel que $\mathfrak S$ dépende algébriquement et intégralement de $\mathfrak R$. Toute base d'intégrité de $\mathfrak R$ est alors complétée en une base d'intégrité de $\mathfrak S$ par une base de $\mathfrak S$ comme module sur un certain sous-anneau de $\mathfrak R$.

Ce critère -- qui présente par lui-même un intérêt autonome -- repose sur l'existence de la base rationnelle et sur le fait que le théorème de la chaîne des diviseurs se conserve lors du passage aux grandeurs entières d'une extension finie. Ce second théorème a été démontré dans la note précédente d'Artin et van der Waerden pour une caractéristique quelconque, mais seulement en posant, pour le domaine initial, une certaine hypothèse de finitude qui est remplie dans le cas des groupes finis (l'anneau des racines représente une extension finie). Pour l'existence de la base rationnelle je donne une démonstration valable pour des corps quelconques comme domaines de coefficients; je peux ainsi démontrer le critère de finitude sous la restriction indiquée ci-dessus.

Que, pour les invariants de groupes finis, le critère de finitude soit rempli repose sur le principe des fonctions symétriques. Tandis qu'en caractéristique nulle les coefficients de la résolvante de Galois fournissent déjà la base d'intégrité, dans le cas général ils n'engendrent que le sous-anneau fini exigé dans le critère de finitude. L'existence de ce sous-anneau fini entraîne, par les mêmes raisonnements, la finitude des invariants ``relatifs'' et des invariants ``modulaires''.

Parmi les invariants considérés ici figurent, comme cas particuliers, les ``invariants projectifs mod $p$'' traités par Dickson et son école.\footnote{Voir Dickson, On Invariants and the theory of Numbers. The Madison Colloquium 1913. La littérature plus récente se trouve dans les volumes des Transactions of the American Mathematical Society parus depuis lors.} Ici le théorème de finitude était connu pour les invariants modulaires, mais non pour les invariants ``formels'' généraux.

\subsection*{\S{} 1. Le critère de finitude}

1. \textbf{Théorème de l'existence de la base rationnelle.} Première formulation: Tout système $S$ de fonctions rationnelles de $n$ indéterminées, à coefficients dans un corps $P$, possède une base rationnelle finie; c'est-à-dire qu'il existe dans le système un nombre fini de fonctions telles que toute fonction du système puisse s'exprimer rationnellement, à coefficients dans $P$, au moyen de ces fonctions en nombre fini.

Seconde formulation: Soit $\mathfrak K$ un corps intermédiaire entre $P$ et $P(x_1,\ldots,x_n)$ -- où $x_1,\ldots,x_n$ désignent des indéterminées -- de degré de transcendance $t\leq n$; soit $y_1,\ldots,y_t$ un système irréductible\footnote{Un système est dit irréductible, d'après Steinitz, lorsqu'il est formé de fonctions algébriquement indépendantes. Un système est dit de degré de transcendance $t$ lorsqu'il contient au moins un système irréductible de $t$ éléments, mais aucun de plus d'éléments. Pour les théorèmes simples sur les systèmes irréductibles utilisés ci-dessous, voir Steinitz, \S{} 22.} de $\mathfrak K$. Alors $\mathfrak K$ devient une extension algébrique finie de $P(y_1,\ldots,y_t)$.

Les deux formulations sont en fait équivalentes. Il suffit en effet de démontrer l'existence de la base rationnelle pour tous les corps intermédiaires; car le corps dérivé d'un système $S$ devient un tel corps intermédiaire $\mathfrak K$, dont la base rationnelle fournit immédiatement une base rationnelle de $S$, puisque chaque élément de $\mathfrak K$ devient une combinaison rationnelle d'un nombre fini de fonctions de $S$. Mais il résulte de la seconde formulation qu'un système irréductible arbitraire $y_1,\ldots,y_t$ de $\mathfrak K$, augmenté des fonctions en nombre fini qui caractérisent $\mathfrak K$ comme extension finie de $P(y_1,\ldots,y_t)$, fournit une telle base rationnelle. Réciproquement, si, d'après la première formulation, une base rationnelle de $\mathfrak K$ est donnée, alors, par la définition du degré de transcendance, chaque élément de base doit être algébriquement dépendant de $P(y_1,\ldots,y_t)$; ainsi $\mathfrak K$ se trouve caractérisé comme extension algébrique finie de $P(y_1,\ldots,y_t)$.

Le théorème sera démontré dans la seconde formulation. Supposons d'abord que $\mathfrak K$ ait le même degré de transcendance $n$ que $P(x_1,\ldots,x_n)$; soit $y_1,\ldots,y_n$ un système irréductible de $\mathfrak K$ et par conséquent de $P(x_1,\ldots,x_n)$. À cause du degré de transcendance $n$ de $y_1,\ldots,y_n$, chaque $x_i$ devient algébrique sur $P(y_1,\ldots,y_n)$; donc $P(x_1,\ldots,x_n)$, et par là $\mathfrak K$ comme corps intermédiaire, devient une extension algébrique finie de $P(y_1,\ldots,y_n)$.

Soit maintenant $t<n$; que la numérotation des $x$ soit choisie de telle sorte que $y_1,\ldots,y_t,x_{t+1},\ldots,x_n$ forment un système irréductible dans $P(x_1,\ldots,x_n)$. Ainsi $x_{t+1},\ldots,x_n$ sont irréductibles relativement à $P(y_1,\ldots,y_t)$ et donc, d'après la loi transitive de la dépendance algébrique, également irréductibles relativement à $\mathfrak K$ et à tout corps intermédiaire $\mathfrak M$ entre $P(y_1,\ldots,y_t)$ et $\mathfrak K$. Cela signifie que $\overline{\mathfrak K}=\mathfrak K(x_{t+1},\ldots,x_n)$, respectivement $\overline{\mathfrak M}=\mathfrak M(x_{t+1},\ldots,x_n)$, devient une extension purement transcendante de $\mathfrak K$, respectivement de $\mathfrak M$; donc tout élément de $\overline{\mathfrak K}$ qui n'appartient pas à $\mathfrak K$ devient transcendant relativement à $\mathfrak K$, et de même tout élément de $\overline{\mathfrak M}$ qui n'appartient pas à $\mathfrak M$ devient transcendant relativement à $\mathfrak M$. De même, posons $\overline P$ égal à l'extension purement transcendante $P(x_{t+1},\ldots,x_n)$; alors $y_1,\ldots,y_t$ sont irréductibles relativement à $\overline P$.

La démonstration va maintenant se ramener au cas déjà traité en prenant $\overline P$ comme domaine des coefficients. De $P<\mathfrak K<P(x_1,\ldots,x_t,x_{t+1},\ldots,x_n)$\footnote{$P<\mathfrak K$ signifie: $P$ est contenu dans $\mathfrak K$, donc est un sous-corps de $\mathfrak K$.} il résulte en effet que $\overline P<\overline{\mathfrak K}<\overline P(x_1,\ldots,x_t)$, où $\overline{\mathfrak K}$ devient, relativement à $\overline P$, du même degré de transcendance $t$ que $\overline P(x_1,\ldots,x_t)$. Par conséquent $\overline{\mathfrak K}$ possède, relativement à $\overline P$ comme domaine des coefficients, une base rationnelle finie, disons $y_1,\ldots,y_t; z_1,\ldots,z_s$. Les $z$ peuvent être choisis dans $\mathfrak K$, puisque $\overline{\mathfrak K}$ est le corps composé de $\mathfrak K$ et de $\overline P$. Il faut montrer que $\mathfrak K$ est égal à $\mathfrak M=P(y_1,\ldots,y_t;z_1,\ldots,z_s)$. Or $\mathfrak M$ est un corps intermédiaire entre $P(y_1,\ldots,y_t)$ et $\mathfrak K$; donc tout élément de $\mathfrak K$ est algébrique relativement à $\mathfrak M$. D'autre part $\overline{\mathfrak M}$ est égal à $\overline{\mathfrak K}$ et ainsi $\mathfrak K$ est contenu dans $\overline{\mathfrak M}$. Donc $\mathfrak K$ devient un corps intermédiaire entre $\mathfrak M$ et $\overline{\mathfrak M}$; et comme tout élément de $\overline{\mathfrak M}$ qui n'appartient pas à $\mathfrak M$ est transcendant relativement à $\mathfrak M$, il s'ensuit nécessairement que $\mathfrak K$ est identique à $\mathfrak M$.

Conséquence: Si $P<\mathfrak K<\mathfrak L<P(x_1,\ldots,x_n)$ et si $\mathfrak L$ est algébrique relativement à $\mathfrak K$, alors $\mathfrak L$ est une extension algébrique finie de $\mathfrak K$. Car chaque élément d'une base rationnelle de $\mathfrak L$ est alors algébrique relativement à $\mathfrak K$; ainsi $\mathfrak L$ est caractérisé comme extension finie de $\mathfrak K$.

2. \textbf{Démonstration du critère de finitude.} La condition est évidemment nécessaire; car si $\mathfrak S$ est un domaine d’intégrité fini, alors $\mathfrak S$ lui-même peut être pris pour le sous-anneau fini $\mathfrak R$ qui intervient dans le critère.

La condition est suffisante sous l'hypothèse que le corps des racines $P^{1/p}$ est une extension finie du corps des coefficients\footnote{Voir p. 28, note 2.} lorsque celui-ci est de caractéristique $p$; pour la caractéristique nulle, aucune hypothèse restrictive n'est nécessaire. Pour le montrer, il faut prouver que $\mathfrak S$ possède une base finie comme module sur un sous-anneau de $\mathfrak R$.

Soit $\mathfrak K$ le corps des fractions\footnote{Voir p. 28, note 2.} de $\mathfrak R$ et soit $\mathfrak L$ le corps des fractions de $\mathfrak S$. Ces corps des fractions existent, puisqu'il s'agit d'anneaux sans diviseurs de zéro; et l'on a $P<\mathfrak K<\mathfrak L<P(x_1,\ldots,x_n)$. Donc, d'après la conséquence sous 1., $\mathfrak L$ devient un corps d'extension algébrique fini de $\mathfrak K$; en outre, par hypothèse, chaque élément de $\mathfrak S$ est entier relativement à $\mathfrak R$, et $\mathfrak S$ devient ainsi un sous-anneau de l'anneau $\mathfrak S'$ formé de tous les éléments de $\mathfrak L$ entiers sur $\mathfrak R$. Mais comme rien n'est supposé sur la clôture intégrale de $\mathfrak R$ dans $\mathfrak K$, le fait que le théorème de la chaîne des diviseurs se conserve lors d'une extension finie ne peut pas être invoqué directement. Il faut plutôt passer d'abord à un sous-anneau $\mathfrak T$ de $\mathfrak R$, lui aussi fini, pour lequel cette clôture intégrale est satisfaite.

Supposons d'abord que le domaine des coefficients $P$ contienne une infinité d'éléments. Alors on peut choisir dans $\mathfrak R$ un système irréductible $g_1(x),\ldots,g_t(x)$ de telle sorte que $\mathfrak R$, et donc aussi $\mathfrak S$, dépende intégralement de $\mathfrak T=P[g_1(x),\ldots,g_t(x)]$. En effet, soit $f_1(x),\ldots,f_k(x)$ la base d'intégrité de $\mathfrak R$ qui existe par hypothèse; si celle-ci ne forme pas un système irréductible, alors -- après une éventuelle renumérotation des indices -- soit $F(f_k;f_1,\ldots,f_{k-1})=0$ une relation à laquelle satisfait $f_k$. Si l'on remplace $f_1,\ldots,f_{k-1}$ par
\[
 f_1+t_1f_k,\ldots,f_{k-1}+t_{k-1}f_k,
\]
alors $f_k$, et donc aussi $f_1,\ldots,f_{k-1}$, dépend intégralement de ces nouvelles fonctions lorsque les $t_i$ sont choisis comme indéterminées; puisque $P$ possède une infinité d'éléments, les $t_i$ peuvent être spécialisés en des éléments de $P$ de telle sorte que cette dépendance intégrale soit conservée. Après un nombre fini d'étapes, en tenant compte de la loi transitive de la dépendance intégrale, on parvient donc à l'anneau cherché $\mathfrak T=P[g_1(x),\ldots,g_t(x)]$\footnote{Voir Hilbert, Über die vollen Invariantensysteme, \S{} 1. Math. Ann. 42 (1893), p. 313--373.}. Le corps $\mathfrak L$ devient une extension finie du corps des fractions $P(g_1(x),\ldots,g_t(x))$ de $\mathfrak T$, et l'anneau $\mathfrak S'$ formé de tous les éléments de $\mathfrak L$ entiers sur $\mathfrak R$ devient identique à celui de tous les éléments de $\mathfrak L$ entiers sur $\mathfrak T$.

Dans $\mathfrak T$ vaut le théorème de la chaîne des diviseurs pour le système de tous les idéaux, et $\mathfrak T$ est intégralement clos dans son corps des fractions, les deux choses parce que $\mathfrak T$ est isomorphe à l'anneau de polynômes en $t$ indéterminées. De l'hypothèse que $P^{1/p}$ est fini relativement à $P$, il résulte en outre que $\mathfrak T^{1/p}$ est fini relativement à $\mathfrak T$.\footnote{Voir la note d'Artin--van der Waerden citée dans l'introduction.} Ainsi vaut dans $\mathfrak S'$ le théorème de la chaîne des diviseurs pour le système de tous les $\mathfrak T$-modules -- en particulier $\mathfrak S$, comme sous-anneau de $\mathfrak S'$ contenant $\mathfrak T$, possède une base finie comme $\mathfrak T$-module\footnotemark[\value{footnote}]. L'existence d'une telle base finie de $\mathfrak S$ comme $\mathfrak T$-module reste encore valable sans aucune hypothèse restrictive sur le domaine des coefficients $P$, lorsque $\mathfrak L$ devient une extension de première espèce de $P(g_1(x),\ldots,g_t(x))$, donc en particulier toujours lorsque $P$ est de caractéristique nulle.\footnote{Voir par exemple E. Noether, Abstrakter Aufbau der Idealtheorie in algebraischen Zahl- und Funktionenkörpern, \S{} 3. Math. Ann. 96 (1926), p. 26--61.}

Soit $h_1(x),\ldots,h_s(x)$ une base de $\mathfrak S$ comme $\mathfrak T$-module; alors la représentation
\[
 f(x)=A_1(g)h_1(x)+\cdots+A_s(g)h_s(x)
\]
pour un $f(x)$ quelconque de $\mathfrak S$ -- où les $A$, comme éléments de $\mathfrak T$, désignent des polynômes en les $g$ -- montre que $g_1(x),\ldots,g_t(x),h_1(x),\ldots,h_s(x)$ forment la base d'intégrité cherchée.

Si le corps $P$ ne contient qu'un nombre fini d'éléments, le corps $P(u)$ obtenu par adjonction d'une indéterminée $u$ -- c'est-à-dire d'un élément transcendant relativement à $\mathfrak S$ -- contient cependant une infinité d'éléments. Ainsi l'anneau $\mathfrak S(u)$ possède une base d'intégrité finie relativement à $P(u)$ comme domaine des coefficients; et comme $\mathfrak S(u)$ est l'anneau composé de $\mathfrak S$ et de $P(u)$, cette base peut être choisie dans $\mathfrak S$ -- par décomposition suivant les puissances de $u$ -- bien que les éléments de base $g_i$ de $\mathfrak S$ ne forment plus maintenant un système irréductible. Ainsi tout polynôme $f(x)$ de $\mathfrak S$ admet une représentation $C(u)f(x)=G(u,g_i,h_i)$, où $C(u)$ désigne un polynôme non nul de $P[u]$ et où $G$ devient entier en $u,g_i,h_i$. La comparaison des coefficients d'une puissance convenable de $u$ montre que les $g_i(x)$ et les $h_i(x)$ donnent une base d'intégrité de $\mathfrak S$. Le \textbf{critère de finitude est ainsi démontré en toute généralité}.

\subsection*{\S{} 2. La finitude des invariants}

1. Par \textbf{groupe linéaire fini de caractéristique $p$}, on entend un groupe $\mathfrak H$ de $h$ substitutions linéaires $A_1,\ldots,A_h$ (de déterminant non nul) à coefficients dans un corps $P$ de caractéristique $p$. Ici $A_k$ représente la substitution linéaire
\[
\begin{aligned}
 x_1^{(k)}&=a_{11}^{(k)}x_1+\cdots+a_{1n}^{(k)}x_n,\\[-0.2em]
 &\ldots,\qquad
 x_n^{(k)}=a_{n1}^{(k)}x_1+\cdots+a_{nn}^{(k)}x_n,
\end{aligned}
\]
ou, en abrégé: $(x^{(k)})=A_k(x)$. Le groupe $\mathfrak H$ transforme donc la suite $(x)$ aux éléments $x_1,\ldots,x_n$ en les suites $(x^{(k)})$ aux éléments $x_1^{(k)},\ldots,x_n^{(k)}$. Comme l'identité doit être contenue parmi $A_1,\ldots,A_h$, la suite $(x)$ est elle aussi contenue parmi les suites $(x^{(k)})$.

Par \textbf{invariant rationnel entier (absolu) du groupe}, on entend une telle fonction rationnelle entière de $x_1,\ldots,x_n$ -- à coefficients dans $P$\footnote{Cela ne constitue pas une restriction de la généralité, dans la mesure où tout invariant à coefficients dans un corps de caractéristique $p$ -- lequel doit posséder avec $P$ un corps commun qui les contienne afin que les invariants soient définis -- peut se décomposer linéairement en un nombre fini d'invariants à coefficients dans $P$. Dans le cas général, il suffit d'exprimer les coefficients par des coefficients linéairement indépendants relativement à $P$; alors la propriété d'invariance est satisfaite identiquement en ces coefficients indépendants.} -- qui reste identiquement inchangée lorsqu'on applique $A_1,\ldots,A_h$. À ces invariants appartiennent donc toutes les fonctions symétriques -- à coefficients dans $P$ -- des suites $(x^{(k)})$. Réciproquement, pour chaque invariant de ce genre, on a
\[
 f(x)=f(x^{(1)})=\cdots=f(x^{(h)})
 \quad\text{ou}\quad
 h\cdot f(x)=f(x^{(1)})+\cdots+f(x^{(h)}).
\]
Il s'ensuit que, dans le cas où l'ordre $h$ du groupe n'est pas divisible par la caractéristique -- donc en particulier dans le cas de la caractéristique nulle -- les fonctions symétriques des suites $(x^{(k)})$, à coefficients dans $P$, épuisent aussi la totalité des invariants. Comme ces fonctions symétriques possèdent une base d'intégrité finie, la démonstration de finitude est achevée dans ce cas; en même temps, la base d'intégrité s'obtient comme le système $G_{\alpha\alpha_1\ldots\alpha_n}(x)$ des coefficients de la résolvante de Galois\footnote{Voir p. 28, note 1.}:
\[
 \Phi(z,u)=\prod_{k=1}^{h}\bigl(z-u_1x_1^{(k)}-\cdots-u_nx_n^{(k)}\bigr)
\]
\[
 =z^h+\sum G_{\alpha\alpha_1\ldots\alpha_n}(x)
 z^{\alpha}u_1^{\alpha_1}\cdots u_n^{\alpha_n}
 \quad(\alpha\ne h;\ \alpha+\alpha_1+\cdots+\alpha_n=h).
\]

2. Dans le cas où $h$ est divisible par $p$, il n'est pas nécessaire que chaque invariant du groupe soit une fonction symétrique des suites $(x^{(k)})$, puisqu'on ne peut plus conclure de $h\cdot f(x)$ à $f(x)$. Mais l'anneau $\mathfrak R$ dérivé des invariants en nombre fini $G_{\alpha\alpha_1\ldots\alpha_n}(x)$ -- les coefficients de la résolvante de Galois -- forme un sous-anneau fini du système $\mathfrak S$ de tous les invariants, et il reste à montrer que $\mathfrak S$ remplit, relativement à $\mathfrak R$, les hypothèses du critère de finitude.

On peut d'abord, sans restreindre la généralité\footnote{Voir p. 33, note.}, prendre comme domaine des coefficients le sous-corps $\overline P$ de $P$ dérivé de tous les coefficients de substitution $a_{\mu\nu}^{(k)}$. Ce corps, dérivé d'un nombre fini d'éléments, naît donc par adjonction au corps premier d'un nombre fini d'éléments algébriques ou transcendants; et, puisque le corps premier est parfait, il s'ensuit que $\overline P^{1/p}$ est fini relativement à $\overline P$.\footnote{Voir p. 32, note 2.}

Il faut encore montrer que $\mathfrak S$ dépend intégralement de $\mathfrak R$. D'après 1., parmi les $h$ suites $x^{(k)}$ apparaît aussi la suite initiale $x_1,\ldots,x_n$; donc la résolvante de Galois $\Phi(z,u)$ s'annule identiquement en $u$ pour $z=u_1x_1+\cdots+u_nx_n$. Si l'on spécialise chaque fois tous les $u$ à zéro sauf $u_i$, qui est posé égal à l'unité, on obtient une relation
\[
 x_i^h+\sum x_i^\alpha G_{\alpha0\ldots\alpha_i\ldots0}(x)=0
 \quad (\alpha\ne h;\ \alpha+\alpha_i=h),
\]
qui montre que chaque $x_i$ dépend intégralement de $\mathfrak R$. Il en va donc de même pour un polynôme quelconque en les $x$; en particulier $\mathfrak S$ dépend intégralement de $\mathfrak R$. Les hypothèses du critère de finitude sont ainsi remplies, et la \textbf{finitude des invariants est démontrée}.

3. Il faut remarquer que les mêmes raisonnements donnent la \textbf{finitude des invariants relatifs et modulaires}. Un polynôme $f(x)$ s'appelle \textbf{invariant relatif} du groupe si tous les $f(x^{(k)})$ ne diffèrent de $f(x)$ que par un facteur non nul de $\overline P$. Le système des invariants absolus -- en particulier donc l'anneau $\mathfrak R$ dérivé des $G_{\alpha\alpha_1\ldots\alpha_n}(x)$ -- forme donc un sous-anneau fini de l'anneau dérivé de tous les invariants relatifs; et, comme ci-dessus, les hypothèses du critère de finitude relativement à $\mathfrak R$ sont remplies.

Un polynôme $f(x)$ s'appelle \textbf{invariant modulaire} du groupe lorsque $f(x)$ devient égal à $f(x^{(k)})$ pour tout système spécial de valeurs $x_1=\alpha_1,\ldots,x_n=\alpha_n$ de $\overline P$. Tout invariant absolu est en même temps invariant modulaire; mais, si $\overline P$ ne contient qu'un nombre fini d'éléments, la réciproque n'a pas besoin d'être vraie. Ici encore les hypothèses du critère de finitude sont remplies relativement au sous-anneau $\mathfrak R$ dérivé des $G_{\alpha\alpha_1\ldots\alpha_n}(x)$.
